\caption{Two content domains without a wall: a hopper's talk follows the room it is said in (H102). (a)~Schematic: the coupling model puts a hopper inside the wall ($s\approx\kappa p$, with $p$ its read share from the other room); observed, its home-room talk stays in its own domain. (b)~Bimodality $D$ of the stayers' positions on the inter-room axis per unit; gray: room-relabel null (axis refitted), 95th percentile. Rooms with different work form two domains ($D$ 5.3--6.1 in \#38, \#41, \#44 and 5.5 in \#51g; relabel $p\le0.006$); identical-kickoff weeks split less (median 1.70). (c)~Unit \#51g: wall coordinate $s$ of the \#general stayers' and \#focus residents' statements, the stayers' leave-one-out positions (rug) and the hoppers' mean positions for statements made in \#general (green) and in \#focus (pink). The one frequent hopper, Claude Fable~5, sits at $s=0.06$ in \#general talk, inside the stayers' range (95th percentile 0.21); hoppers' statements made inside \#focus sit at $s$ 0.15--0.51. (d)~Hopper-day $s$ of \#general talk against hopping read-outs (within-hopper): $\kappa_R=-0.014$ [$-0.042$, $0.024$] per e-fold (band: day-bootstrap 95\%), where the synthetic power against $\kappa=0.1$ (dashed) is $\ge0.93$. bge, style-residualized statement vectors.}
\label{fig:H102}
